\documentclass[book]{jlreq}
\usepackage{luatexja}
\usepackage{subfiles}
\usepackage{amsmath,amssymb,bm}
\usepackage{url}
\usepackage[nottoc,notlof,notlot,notindex]{tocbibind}
\urlstyle{same}

% 文献番号と本文の間隔を確保する。
\makeatletter
\renewcommand{\@openbib@code}{%
  \setlength{\labelsep}{0.75em}%
  \setlength{\leftmargin}{\labelwidth}%
  \addtolength{\leftmargin}{\labelsep}%
}
\makeatother

% 章別PDFでは、その章で引用した文献を末尾に出力する。
\newcommand{\chapterbibliography}{%
  \ifSubfilesClassLoaded{%
    \bibliographystyle{unsrt}%
    \bibliography{../references}%
  }{}%
}

\title{ニューラルネットワークの最適化\\[0.5em]
  {\large 勾配の前処理から考える}}
\author{石川智貴}
\date{}

\begin{document}
\maketitle

\tableofcontents

% 各ファイルは単独でもコンパイルできます。
\subfile{sec/01-introduction}
\subfile{sec/02-mirror-descent}
\subfile{sec/03-constrained-gradient}
\subfile{sec/04-minibatch-learning}
\subfile{sec/05-linear-networks}
\subfile{sec/06-infinite-width}
\subfile{sec/07-hopfield}
\subfile{sec/08-hyperparameters}

\appendix
\subfile{sec/appendix-ntk-stability}

\bibliographystyle{unsrt}
\bibliography{references}

\end{document}
